\section{Detailed Description of FL-RL-TRC}

In this section, the FL-RL-TRC caching mechanism is explained in detail, covering the two-phase architecture that combines Federated Learning for proactive cache pre-placement with Reinforcement Learning for reactive cache eviction.

\subsection{Task Popularity Prediction using Federated Learning}

Efficient cache management in VANETs requires predicting which task results will be requested in the future to enable proactive pre-placement. Constantly querying all RSUs about their cache states would create significant network overhead and consume bandwidth needed for critical vehicular communications. This section proposes a Federated Learning-based predictive model for estimating task popularity patterns across distributed fog nodes, thereby minimizing communication overhead while enabling intelligent cache pre-placement decisions.

The FL-based prediction process begins with feature engineering and distributed data collection across RSUs, followed by federated model training that preserves data privacy. Finally, the trained model predicts task popularity scores to guide cache pre-placement decisions.

\subsubsection{Feature Engineering for Task Popularity}

To accurately predict task popularity, we developed a comprehensive 12-dimensional feature vector that captures multiple aspects of task characteristics, network conditions, and temporal dynamics. The feature engineering process utilizes task request history collected at each RSU over time. Table \ref{table:fl_feature_descriptions} details the features used in the FL advisor.

\begin{table}[t!]
    \centering
    \caption{Feature descriptions for FL-based task popularity prediction}
    \begin{footnotesize}
    \begin{tabularx}{\linewidth}{>{\centering\arraybackslash}m{2.5cm}|>{\RaggedRight\arraybackslash}X}
    \hline
    \rowcolor[gray]{0.9}
    \textbf{Feature Name} & \textbf{Description} \\
    \hline
    RequestFreq & Normalized request frequency: number of times this task type has been requested divided by total requests \\
    \hline
    LogFreq & Log-scaled frequency: $\log(1 + \text{count})$ to handle skewed distributions \\
    \hline
    Priority & Task priority level (0-3), with higher values indicating greater urgency \\
    \hline
    DataSize & Task data payload size in kilobytes, normalized to [0,1] range \\
    \hline
    IsLocal & Binary indicator: 1 if task originated in this RSU's coverage area, 0 otherwise \\
    \hline
    VehicleSpeed & Average speed of vehicles requesting this task type, normalized \\
    \hline
    RSULoad & Current computational load at the RSU, measured as percentage of capacity utilized \\
    \hline
    PopularityDecay & Exponential decay factor: $e^{-\lambda \cdot \Delta t}$ where $\Delta t$ is time since last request and $\lambda = 0.2$ \\
    \hline
    TaskType\_0 & One-hot encoding for task type 0 (e.g., route optimization) \\
    \hline
    TaskType\_1 & One-hot encoding for task type 1 (e.g., traffic analysis) \\
    \hline
    TaskType\_2 & One-hot encoding for task type 2 (e.g., sensor processing) \\
    \hline
    TaskType\_3 & One-hot encoding for task type 3 (e.g., hazard detection) \\
    \hline
    \end{tabularx}
    \end{footnotesize}
    \label{table:fl_feature_descriptions}
\end{table}

Data preparation involves collecting task request histories at each RSU, including task characteristics, request timestamps, and execution outcomes. Feature normalization ensures all features are scaled to similar ranges, preventing dominance by large-value features. The popularity decay parameter $\lambda = 0.2$ was determined through systematic experimentation across values \{0.01, 0.02, 0.05, 0.1, 0.2\}, with $\lambda = 0.2$ providing optimal performance by capturing the rapid dynamics of vehicular environments where task relevance diminishes with approximately 5-second half-life.

\subsubsection{Federated Learning Model Architecture}

The FL advisor employs a Multi-Layer Perceptron (MLP) neural network architecture for task popularity prediction. The model architecture consists of:

\begin{itemize}
    \item \textbf{Input Layer:} 12 neurons corresponding to the 12-dimensional feature vector
    \item \textbf{Hidden Layer 1:} 64 neurons with ReLU activation
    \item \textbf{Hidden Layer 2:} 32 neurons with ReLU activation
    \item \textbf{Output Layer:} 1 neuron with sigmoid activation, producing popularity score in [0,1] range
\end{itemize}

The federated learning process follows the Federated Averaging (FedAvg) algorithm \cite{mcmahan2017communication}:

\begin{enumerate}
    \item \textbf{Initialization:} The SDN controller initializes global model parameters $\theta_0$
    \item \textbf{Local Training:} Each RSU $m$ trains the model on its local dataset $D_m$ for $E$ epochs:
    \begin{equation}
        \theta_m^{t+1} = \theta_m^t - \eta \nabla L(\theta_m^t; D_m)
    \end{equation}
    where $\eta$ is the learning rate and $L$ is the binary cross-entropy loss function
    \item \textbf{Model Aggregation:} The controller aggregates local models using weighted averaging:
    \begin{equation}
        \theta^{t+1} = \sum_{m=1}^{M} \frac{|D_m|}{|D|} \theta_m^{t+1}
    \end{equation}
    where $|D_m|$ is the size of local dataset at RSU $m$ and $|D| = \sum_{m=1}^{M} |D_m|$
    \item \textbf{Distribution:} The updated global model is distributed back to all RSUs
\end{enumerate}

This federated approach enables collaborative learning across RSUs while preserving data privacy, as raw task data never leaves individual fog nodes.

\subsubsection{Model Evaluation and Performance}

The FL advisor's performance was evaluated using standard classification metrics. The model predicts a binary label indicating whether a task result should be pre-placed in cache (1) or not (0), based on a threshold of 0.5 applied to the popularity score. Table \ref{table:fl_model_performance} summarizes the FL advisor's performance across different configurations.

\begin{table}[t]
\centering
\caption{FL advisor performance with different popularity decay parameters}
\begin{tabular}{c|c|c|c}
\hline
\rowcolor[gray]{0.9}
\textbf{$\lambda$} & \textbf{Precision (\%)} & \textbf{Recall (\%)} & \textbf{F1-Score (\%)} \\
\hline
0.01 & 72.3 & 68.5 & 70.3 \\
\hline
0.05 & 78.1 & 74.2 & 76.1 \\
\hline
0.1 & 82.4 & 79.8 & 81.1 \\
\hline
0.2 & 86.7 & 84.3 & 85.5 \\
\hline
\end{tabular}
\label{table:fl_model_performance}
\end{table}

The evaluation revealed that $\lambda = 0.2$ provides superior performance, achieving 86.7\% precision and 84.3\% recall. Fig. \ref{fig:fl_training_curve} illustrates the convergence of the FL model over federated training rounds, demonstrating stable learning across distributed RSUs.

\begin{figure}[t!]
\centering
\includegraphics[width=0.75\linewidth]{fl_training_curve.pdf}
\caption{FL advisor training convergence across federated rounds.}
\label{fig:fl_training_curve}
\end{figure}

\subsection{Dynamic Cache Eviction using Reinforcement Learning}

While the FL advisor handles proactive cache pre-placement, real-time cache management during task execution requires adaptive eviction policies that respond to dynamic network conditions. The RL agent continuously monitors cache states and makes intelligent eviction decisions when storage capacity is reached, balancing multiple competing objectives such as maximizing cache hit rates, minimizing task completion times, and ensuring deadline compliance.

\subsubsection{State Space Representation}

The RL agent observes a 9-dimensional state space that captures the current cache status, network conditions, and task characteristics. The state vector $s_t$ at timestep $t$ is defined as:

\begin{equation}
s_t = [u_t, h_t, c_t, \bar{p}_t, \bar{s}_t, \bar{d}_t, l_t, q_t, a_t]
\end{equation}

where:
\begin{itemize}
    \item $u_t$: Cache usage ratio (occupied slots / total capacity)
    \item $h_t$: Current cache hit rate over recent window
    \item $c_t$: Average task completion time (normalized)
    \item $\bar{p}_t$: Mean priority of cached tasks
    \item $\bar{s}_t$: Mean data size of cached tasks (normalized)
    \item $\bar{d}_t$: Mean remaining TTL of cached entries (normalized)
    \item $l_t$: RSU computational load (percentage)
    \item $q_t$: Task queue length (number of pending tasks)
    \item $a_t$: Age of oldest cache entry (normalized)
\end{itemize}

All continuous features are normalized to [0,1] range to facilitate neural network training.

\subsubsection{Action Space Definition}

The RL agent selects from five discrete actions at each decision point:

\begin{equation}
\mathcal{A} = \{a_0, a_1, a_2, a_3, a_4\}
\end{equation}

where:
\begin{itemize}
    \item $a_0$: \textbf{No Operation} - Maintain current cache state
    \item $a_1$: \textbf{Evict LRU} - Remove least recently used entry
    \item $a_2$: \textbf{Evict Random} - Remove random entry (exploration)
    \item $a_3$: \textbf{Evict LFU} - Remove least frequently used entry
    \item $a_4$: \textbf{Shrink TTL} - Halve TTL of all entries (aging policy)
\end{itemize}

This action space provides flexibility for the agent to learn context-dependent eviction strategies rather than relying on a single fixed policy.

\subsubsection{Reward Function Design}

The reward function balances multiple objectives to guide the RL agent toward optimal caching decisions. At each timestep $t$, the reward $r_t$ is computed as:

\begin{equation}
r_t = w_h \cdot r_h + w_s \cdot r_s + w_c \cdot r_c + w_e \cdot r_e
\end{equation}

where the components are:

\begin{itemize}
    \item \textbf{Cache Hit Reward} ($r_h$): 
    \begin{equation}
        r_h = \begin{cases}
            +3.0, & \text{if cache hit} \\
            -1.0, & \text{if cache miss}
        \end{cases}
    \end{equation}
    
    \item \textbf{Success Rate Reward} ($r_s$):
    \begin{equation}
        r_s = \begin{cases}
            +2.0, & \text{if task completed before deadline} \\
            -2.0, & \text{if task missed deadline}
        \end{cases}
    \end{equation}
    
    \item \textbf{Completion Time Reward} ($r_c$):
    \begin{equation}
        r_c = -\frac{t_{\text{completion}}}{t_{\text{deadline}}}
    \end{equation}
    
    \item \textbf{Eviction Penalty} ($r_e$):
    \begin{equation}
        r_e = \begin{cases}
            -0.5, & \text{if evicted high-priority task} \\
            0.0, & \text{otherwise}
        \end{cases}
    \end{equation}
\end{itemize}

The weights are set as $w_h = 3.0$, $w_s = 2.0$, $w_c = 1.0$, $w_e = 1.0$, emphasizing cache hit rate optimization while maintaining task success guarantees.

\subsubsection{Deep Q-Network Architecture}

The RL agent employs a Deep Q-Network (DQN) \cite{mnih2015human} to approximate the action-value function $Q(s, a; \theta)$. The DQN architecture consists of:

\begin{itemize}
    \item \textbf{Input Layer:} 9 neurons (state space dimensions)
    \item \textbf{Hidden Layer 1:} 128 neurons with ReLU activation
    \item \textbf{Hidden Layer 2:} 64 neurons with ReLU activation
    \item \textbf{Output Layer:} 5 neurons (Q-values for each action)
\end{itemize}

The agent uses experience replay with buffer size 10,000 and $\epsilon$-greedy exploration with decay schedule:

\begin{equation}
\epsilon_t = \max(\epsilon_{\min}, \epsilon_{\max} \cdot \gamma^t)
\end{equation}

where $\epsilon_{\max} = 0.3$, $\epsilon_{\min} = 0.01$, and $\gamma = 0.995$.

The Q-learning update rule is:

\begin{equation}
Q(s_t, a_t) \leftarrow Q(s_t, a_t) + \alpha [r_t + \gamma \max_{a'} Q(s_{t+1}, a') - Q(s_t, a_t)]
\end{equation}

where $\alpha = 0.001$ is the learning rate.

\subsubsection{Training and Convergence}

The RL agent is trained online through continuous interaction with the caching environment. Fig. \ref{fig:rl_learning_curve} shows the learning progression, with cache hit rate improving from 87.8\% at 50 timesteps to 97.0\% at 200 timesteps, demonstrating clear learning and convergence.

\begin{figure}[t!]
\centering
\includegraphics[width=0.75\linewidth]{rl_learning_curve.pdf}
\caption{RL agent learning curve showing cache hit rate improvement over training timesteps.}
\label{fig:rl_learning_curve}
\end{figure}

\subsection{Integration of FL and RL Components}

The FL-RL-TRC framework integrates both learning components within the SDN controller architecture. Algorithm \ref{algo:fl_rl_integration} describes the overall caching workflow.

\begin{algorithm}[t!]
\caption{FL-RL-TRC Integrated Caching Framework}
\label{algo:fl_rl_integration}

\SetKwInOut{Input}{Input}
\SetKwInOut{Output}{Output}

\Input{Task stream $\mathcal{T}$, RSU set $\mathcal{M}$, Cache capacity $C$}
\Output{Cache hit rate, Task success rate, Avg completion time}

\Begin{
    Initialize FL advisor model $\theta_{FL}$ via federated training\;
    Initialize RL agent DQN $\theta_{RL}$ with random weights\;
    Initialize experience replay buffer $\mathcal{D}$\;
    
    \tcp{Phase 1: FL-based Pre-placement}
    \For{each RSU $m \in \mathcal{M}$}{
        Collect task history features $X_m$\;
        Predict popularity scores: $P_m = FL\_Advisor(X_m; \theta_{FL})$\;
        Pre-place top-$k$ tasks with highest scores in cache\;
    }
    
    \tcp{Phase 2: RL-based Dynamic Eviction}
    \For{each timestep $t = 1, 2, \ldots, T$}{
        \For{each incoming task $\tau_t$}{
            Check cache for matching result using hash key\;
            
            \If{cache hit}{
                Return cached result immediately\;
                Update hit rate statistics\;
            }
            \Else{
                Execute task at fog node\;
                
                \If{cache is full}{
                    Observe state $s_t$ from cache status\;
                    Select action $a_t$ using $\epsilon$-greedy policy\;
                    Execute eviction action $a_t$\;
                    Store new result in cache\;
                    Observe reward $r_t$ and next state $s_{t+1}$\;
                    Store transition $(s_t, a_t, r_t, s_{t+1})$ in $\mathcal{D}$\;
                    
                    \tcp{DQN Training}
                    Sample mini-batch from $\mathcal{D}$\;
                    Update $\theta_{RL}$ using Q-learning\;
                }
                \Else{
                    Store result directly in cache\;
                }
            }
        }
        
        Decay $\epsilon$ for exploration schedule\;
    }
    
    \tcp{Periodic FL Model Update}
    \If{$t \mod T_{FL} == 0$}{
        Perform federated averaging across RSUs\;
        Update global FL model $\theta_{FL}$\;
        Redistribute to all RSUs\;
    }
}
\end{algorithm}

The integration ensures that FL handles long-term popularity patterns through periodic model updates (e.g., every 100 timesteps), while RL continuously adapts to short-term dynamics through online learning at every decision point.

\subsection{Cache Key Generation and Result Matching}

To enable cache reuse across different vehicles requesting identical computations, FL-RL-TRC uses a hash-based cache key generation mechanism that depends only on task intrinsics, not vehicle identity:

\begin{equation}
key(\tau_{ij}) = \text{hash}(\text{type}_{ij}, \text{priority}_{ij}, \lfloor \text{size}_{ij}/100 \rfloor)
\end{equation}

This approach allows multiple vehicles to benefit from cached results of previously executed identical tasks, maximizing cache utility and reducing redundant computations. The hash function ensures fast lookup with $O(1)$ complexity.

\subsection{Computational Complexity Analysis}

The computational complexity of FL-RL-TRC components is analyzed as follows:

\begin{itemize}
    \item \textbf{FL Advisor Inference:} $O(d \cdot h_1 + h_1 \cdot h_2 + h_2)$ where $d=12$ is input dimension, $h_1=64$, $h_2=32$ are hidden layer sizes. This is negligible compared to task execution time.
    
    \item \textbf{RL Agent Decision:} $O(s \cdot h_1 + h_1 \cdot h_2 + h_2 \cdot a)$ where $s=9$ is state dimension, $a=5$ is action space size. Also negligible.
    
    \item \textbf{Cache Lookup:} $O(1)$ using hash-based indexing.
    
    \item \textbf{Federated Aggregation:} $O(M \cdot |\theta|)$ where $M$ is number of RSUs and $|\theta|$ is model parameter count. Performed periodically, not per-task.
\end{itemize}

The overhead introduced by FL-RL-TRC is minimal compared to task execution times, making it practical for real-time VANET deployments.
